\section{Target partitions, cost learning, and exact checks}
\label{app:batch-policy}
\subsection{Why target grouping changes the seed region}
Let $\mathcal T$ be a finite batch of distinct index vectors from one family.
For a vector $\boldsymbol a$, define its extra denominator powers and numerator
rank as $\delta(\boldsymbol a)=\sum_i\max(a_i-1,0)$ and
$\rho(\boldsymbol a)=\sum_i\max(-a_i,0)$, as in
Appendix~\ref{app:equation-refinement}. A partition
$\Pi=\{\mathcal T_1,\ldots,\mathcal T_K\}$ consists of disjoint nonempty
subsets whose union is $\mathcal T$. The number $K$ is the number of separate
reductions. For each part, the corresponding maxima obey
\begin{equation}
 \max_{\boldsymbol a\in\mathcal T_k}\delta(\boldsymbol a)
 \leq\max_{\boldsymbol a\in\mathcal T}\delta(\boldsymbol a),\qquad
 \max_{\boldsymbol a\in\mathcal T_k}\rho(\boldsymbol a)
 \leq\max_{\boldsymbol a\in\mathcal T}\rho(\boldsymbol a).
 \label{eq:batch-maxima}
\end{equation}
These inequalities follow because each maximum on the left runs over a
subset of the indices on the right. They explain how a partition can lower
one or both seed limits. They do not imply a time saving: separate parts can
generate the same intermediate equations and repeat the preparation.
For example, one part can contain targets with large $\delta$ and zero
$\rho$, while another contains targets with zero $\delta$ and large $\rho$.
Their union requires both large bounds simultaneously. Kira's lower-sector
truncation remains active in every comparison where the fixed seed policy
requests it~\cite{Kira3}.

The five candidate partitions keep the batch together, group equal values of
$\rho$, group equal values of $\delta$, group equal sets of positive indices,
or divide the given target sequence into consecutive groups of at most three.
The first three group definitions depend on integral powers, and the fourth
uses the set of positive indices that defines a sector. The last definition
also depends on the order supplied by the caller. Every target occurs in
exactly one part. The learned method selects one partition before equation
generation and processes its parts sequentially with one Kira worker.
An additional candidate keeps the batch together and disables the equation
filter. Thus the cost predictor compares six calculations. This candidate
allows a target-dependent decision to avoid a filter whose retry cost exceeds
its saving.
When a batch has at most three targets and all targets have equal rank,
equal extra denominator powers, and the same sector, all five partitions
contain just that batch. If the earlier policy also disables its equation
filter on this batch, the application skips loading the cost model and uses
the undivided unfiltered request. This follows directly from the partition
definitions and preserves the available reduction choice.

If a completed part gives
$I_{\boldsymbol a}=\sum_m r_{\boldsymbol a m}M_m$ for every
$\boldsymbol a\in\mathcal T_k$, then the union of these relations supplies
the reduction of every target in $\mathcal T$. Indeed, disjointness assigns
each target to one part, and the union property ensures that no target is
missing. All parts use the same prescribed master basis. For a target
expression $T=\sum_{\boldsymbol a\in\mathcal T}c_{\boldsymbol a}
I_{\boldsymbol a}$, substitution and interchange of the finite sums gives
\begin{equation}
 T=\sum_{k=1}^{K}\sum_{\boldsymbol a\in\mathcal T_k}
       c_{\boldsymbol a}\sum_m r_{\boldsymbol a m}M_m
  =\sum_m\left(\sum_{k=1}^{K}\sum_{\boldsymbol a\in\mathcal T_k}
       c_{\boldsymbol a}r_{\boldsymbol a m}\right)M_m.
 \label{eq:batch-merge}
\end{equation}
The model changes the division of the calculation. Equation~\eqref{eq:batch-merge}
shows how the exact part reductions provide the same requested expression.

\subsection{Physical requests and reference construction}
The denominator momenta, masses, and rational kinematic profiles are those
specified in Appendix~\ref{app:learned-bounds}. The ten structures are the
sunrise, planar double box, massive three-loop vacuum, massive three-loop
ladder, nonplanar box, pentagon, external-virtuality box, three- and four-loop
boxes with bubbles, and four-loop massless two-point family. The dimension
$d$ remains symbolic. For the first seven structures, let $D$, $R$, $X$, and
$P$ denote the three-target sets with extra denominator powers, rank two,
mixed powers, and a pinched propagator in Eq.~\eqref{eq:bounds-targets}.
The requested batches are $D$, $R$, $X$, $D\cup R$, $P\cup R$, and
$D\cup R\cup X\cup P$. The two bubble-box structures and the massless
two-point family use the corresponding one-dot, one-numerator, and corner
sets from that appendix, with the corner set also taking the place of $P$.
Repeated index vectors and identical complete batches at a profile appear
once. Development uses $h=701$, training uses $h=711$, validation uses
$h=712$, and the reserved comparison uses $h=721,722$.

Each component first receives a native reduction with expanded seed limits.
We collect the resulting master indices into a common basis, reduce the
components in that basis, and check each coefficient again after increasing
the numerator limit by one. The checked component relations form the
reference for each union. A failure in discovery, common-basis reduction, or
the enlarged check excludes that batch from the timed comparison and remains
in the record. Agreement with this finite enlargement does not prove a
minimal basis. Reference construction remains outside the application timer;
the timed request receives the basis and generates its own IBP equations.

\subsection{Predictor inputs and objective}
The predictor uses the 27 family and target quantities specified in
Appendix~\ref{app:learned-bounds}, together with five descriptions of the
distribution of powers within the batch:
\begin{equation}
 |\{\delta(\boldsymbol a):\boldsymbol a\in\mathcal T\}|,\quad
 |\{\rho(\boldsymbol a):\boldsymbol a\in\mathcal T\}|,\quad
 \delta_{\max}-\delta_{\min},\quad
 \rho_{\max}-\rho_{\min},\quad
 \frac{1}{|\mathcal T|}\sum_{\boldsymbol a\in\mathcal T}
             \delta(\boldsymbol a)\rho(\boldsymbol a).
\end{equation}
Here $|\cdot|$ denotes the number of members, and the maxima and minima run
over the batch. The last quantity distinguishes targets that carry both
kinds of power from a batch in which different targets carry them separately.
No reference coefficient, measured test time, or kinematic profile number
enters the features.

Candidate definitions can produce identical ordered parts on a particular
batch. Such candidates share the same measured training or validation cost,
and the record identifies the common measurement. The reserved comparison
still executes every complete application method separately, including its
own selection and loading costs.

For training batch $b$ and partition $a$, let $\bar T_{ba}$ be the median of
two complete timings when both reductions succeed. A failed candidate
receives a training cost equal to the larger of 300 seconds and the sum of
its measured attempt times. This chosen penalty makes a bounded failure
costly to the predictor; it is not an observed completion time. With
$a=0$ denoting the undivided batch, the labels and normalized weights are
\begin{equation}
 y_{ba}=\log\frac{\bar T_{ba}}{\bar T_{b0}},\qquad
 u_b=\sqrt{\min_a\bar T_{ba}},\qquad
 w_b=\frac{u_b}{B^{-1}\sum_{b'=1}^{B}u_{b'}}.
\end{equation}
The positive integer $B$ is the number of usable training batches. The
normalization gives $B^{-1}\sum_b w_b=1$ by direct substitution. The fitted
multi-output regressors use weighted squared error,
\begin{equation}
 \mathcal L(\theta)=\sum_{b=1}^{B}w_b\sum_a
       \bigl[f_{\theta,a}(x_b)-y_{ba}\bigr]^2.
 \label{eq:batch-loss}
\end{equation}
Here $x_b$ is the feature vector, $\theta$ denotes the fitted regressor
parameters, and $f_{\theta,a}$ is the predicted log-cost ratio for partition $a$. Since
$y_{ba}-y_{ba'}=\log(\bar T_{ba}/\bar T_{ba'})$ and the logarithm is increasing,
the smallest exact label corresponds to the smallest assigned cost. Fitting
Eq.~\eqref{eq:batch-loss} approximates these labels; it supplies no guarantee
that a prediction is optimal on a new batch.

We fit eighteen candidates: regression trees, random forests with 32 trees,
and extremely randomized forests with 32 trees, each with maximum depth
2, 4, or 6 and minimum leaf population 1 or 3~\cite{CART1984,Breiman2001Forest,Geurts2006ExtraTrees}.
The implementation uses Ref.~\cite{Pedregosa2011Bounds}. Three complete
validation repetitions select a fixed partition $c_g$ for each graph $g$.
A model proposes $a_\star=\arg\min_a f_{\theta,a}(x)$ when
\begin{equation}
 f_{\theta,a_\star}(x)<f_{\theta,c_g}(x)-\gamma,
 \qquad \gamma\in\{0,0.02,0.05,0.1\};
\end{equation}
otherwise it uses $c_g$. Validation chooses the model and margin together.
Deployment requires at least a three-percent saving in summed validation
cost relative to the graph-specific fixed partitions. A separate validation
run with every ML model disabled selects the fixed control for that condition.
Ordered physical structures and target-count and power ranges constrain
deployment; an unmatched input retains an undivided batch.

To measure the remaining opportunity within this candidate pool, we also
sum the smallest recorded cost for each validation batch:
\begin{equation}
 T_{\mathrm{min,recorded}}=\sum_b\min_a\bar T_{ba}
 \leq\sum_b\bar T_{b,c_{g(b)}}=T_{\mathrm{fixed}}.
\end{equation}
Here $g(b)$ is the graph of batch $b$, and the sums run over validation
batches. Each minimum is no greater than the chosen fixed candidate's
cost, which proves the inequality after summation. This retrospective
quantity uses measured costs that a new request does not supply. It is
neither an achieved prediction time nor a lower bound on all possible
reduction methods.

\subsection{Timed scope and failures}
The application loads the policy for each request. Each part uses the fixed
seed configuration for its physical graph. The previous equation filter is
available only inside its original checked domain, with its weights and
threshold unchanged. An insufficient filtered system retries without the
filter, then uses unshifted and expanded seed limits where distinct. A failed
partition can retry the undivided unfiltered calculation if time remains.
The common request budget is 150 seconds, and each native attempt has at most
60 seconds. A request that exhausts this budget remains a failure.

The reserved test rotates five methods through five repetition positions:
the previous undivided workflow, that workflow with every ML model disabled,
the new graph-specific fixed partitions with filters available, independently
tuned fixed partitions with every ML model disabled, and the learned
partition choice with filters available. Every method uses the same exact
implementation, worker count, basis, and coefficient checks. Initial Python
imports precede these repeated timers. The separate application check starts
a new Python process and includes imports, output, and the exact comparison.

For deployment, we retain the frozen domain limits, fixed choices, and
prediction margin. We retain the tree arrays only when validation enables
the predictor. When the predictor is disabled, its choice rule returns
before consulting those arrays, so removing them leaves each choice
unchanged. We check equality of the returned choices on every training,
validation, and reserved input. This check uses no reserved timing or
coefficient outcome. The repeated test retains its original frozen
implementation; the fresh-process application check uses the compact
deployment form. All ML-disabled controls keep their existing behavior.

We sum per-batch median times over requests for which all methods and all
repetitions succeed and whose output is not entirely the identity. Failed
references, incomplete timed batches, and identity-only batches remain
separate. The ratio of two such sums compares the same completed workload.
Paired resampling of the five repetition positions within each batch gives
5,000 ratios; the 2.5th and 97.5th percentiles describe timing variation for
this specified collection. They do not establish transfer to untested
topologies or loop orders. Some index sets recur across rational profiles,
and the report records this overlap.

Independent graph groups run on two workers during data collection and
four workers during the reserved comparison. This shortens
the experiment without parallelizing an individual timed reduction.
Every reserved batch keeps all methods and repetitions on the same worker,
with the same rotated order described above. Complete training measurements
made before this scheduling change remain in the training record.

If a cluster allocation ends during a comparison, we retain each complete
batch and repeat every method and repetition for an unfinished batch on
one replacement worker. The saved partial comparison remains in a separate
archive and does not enter the timing ratios. We assign the remaining whole
batches from validation costs and keep the frozen models, global method
rotation, exact checks, and timing definitions unchanged. The audit checks
the worker identity for every completed comparison batch.
